\chapter{Préparation des données et modélisation du risque de crédit}

Ce chapitre présente les données mobilisées pour ce projet et leur transformation progressive en modèles de risque exploitables. Il part de la source unique dont dispose le projet, en expose la compréhension, le nettoyage et la construction de la variable cible, avant de détailler le développement des modèles de probabilité de défaut puis de perte en cas de défaut. Il se referme sur l'évaluation, la sélection et l'interprétabilité de ces modèles, dont les artefacts entraînés sont ensuite exploités pour la quantification du risque présentée au chapitre suivant.

\section{Compréhension et exploration des données}

\subsection{Présentation de la source des données}

Les données utilisées dans ce projet proviennent d'un unique fichier source réel, \texttt{Base\_depot.xlsx}, extrait du système d'information bancaire. Ce fichier contient des comptes de dépôt, et non des données de crédit : il sert de point d'ancrage, en fournissant une structure de clientèle réelle, clients, comptes, agences et encours, à partir de laquelle les données nécessaires à l'étude du risque de crédit ont ensuite été construites.

\subsection{Description du jeu de données}

Le tableau~\ref{tab:donnees-source} résume les principales caractéristiques de ce jeu de données, avant et après nettoyage.

\begin{table}[htbp]
    \centering
    \footnotesize
    \caption{Chiffres de référence du jeu de données source.}
    \label{tab:donnees-source}
    \begin{tabular}{p{7cm}p{4cm}}
        \toprule
        \textbf{Indicateur} & \textbf{Valeur} \\
        \midrule
        Lignes brutes & 211 825 \\
        Lignes après nettoyage & 211 824 \\
        Clients uniques & 154 619 \\
        Comptes uniques & 211 824 \\
        Agences uniques & 50 \\
        Catégorie à encours élevé (assimilée au segment Entreprises dans la simulation) & 556 clients (0,36\,\%), représentant 40,1\,\% de l'encours total \\
        Catégorie à encours standard (assimilée au segment Particuliers dans la simulation) & 154 063 clients (99,64\,\%), représentant 59,9\,\% de l'encours total \\
        \bottomrule
    \end{tabular}
\end{table}

Cette répartition ne correspond pas à une classification juridique des clients présente dans la source : elle a été construite à partir de leur niveau d'encours afin de définir deux catégories utilisées dans la simulation.

\subsection{Limites de la source et recours aux données synthétiques}

Cette source présente une limite déterminante pour la suite du projet : elle ne contient ni historique de crédit, ni défaut observé, seulement des comptes de dépôt. Les informations relatives aux crédits eux-mêmes, octroi, remboursement, défaut, n'existent donc pas dans les données réelles : elles ont été générées synthétiquement, à partir de la structure réelle de la clientèle et des comptes présentée ci-dessus. Les modèles prédictifs présentés dans la suite de ce chapitre n'ont ainsi pas été entraînés sur des défauts bancaires historiquement observés, mais sur des historiques de défaut générés et calibrés dans ce cadre, dont la construction est détaillée en section~3.3.

\subsection{Traitement des encours négatifs}

Parmi les comptes de la source réelle, 850 présentent un solde d'encours négatif. Ces comptes ont été conservés, plutôt qu'exclus ou corrigés, et signalés par un indicateur dédié : un encours négatif peut constituer une information comportementale potentiellement utile, qu'il aurait été dommageable de supprimer. Cette conservation ne préjuge toutefois pas de la nature de ces encours négatifs, qui ne correspondent pas nécessairement à un incident de paiement.

\section{Nettoyage, qualité et préparation des données}

\subsection{Nettoyage du fichier source}

Le fichier source présente plusieurs anomalies traitées lors de sa préparation. Neuf colonnes présentent un taux de remplissage réel de 0~\% et sont supprimées : le choix a été de ne pas les imputer, aucun contenu réel ne permettant de justifier une valeur imputée. Une ligne contient par ailleurs un en-tête de colonne inséré accidentellement lors d'une exportation ; elle est supprimée par un filtre sur la devise. C'est cette même ligne qui présente l'unique valeur d'encours manquante du fichier, également supprimée par ce filtre. Les encours négatifs, traités séparément, ont été présentés en section~3.1.

\subsection{Construction de la variable de segmentation}

Pour distinguer les clients des segments Particuliers et Entreprises, introduits en section~3.1, une variable d'encours agrégé par client est d'abord calculée, en sommant l'encours de l'ensemble des comptes qu'il détient, y compris ses éventuels encours négatifs. Un client est affecté au segment Entreprises (simulé) si cet encours agrégé atteint ou dépasse 50 millions de XOF, et au segment Particuliers (simulé) en deçà de ce seuil.

Ce seuil a été retenu après une analyse de sensibilité portant sur huit valeurs candidates, dont les résultats sont résumés dans le tableau~\ref{tab:seuil-segmentation}. Pour chacune, l'écart entre l'encours minimum du segment Entreprises et l'encours maximum du segment Particuliers a été mesuré.

\begin{table}[htbp]
    \centering
    \footnotesize
    \caption{Analyse de sensibilité du seuil de segmentation.}
    \label{tab:seuil-segmentation}
    \begin{tabular}{p{2.2cm}p{2.8cm}p{1.8cm}p{1.8cm}p{2.8cm}}
        \toprule
        \textbf{Seuil} & \textbf{Clients segment élevé} & \textbf{\% clients} & \textbf{\% encours} & \textbf{Écart de séparation} \\
        \midrule
        10 M XOF & 4 789 & 3,10\,\% & 68,4\,\% & 100 XOF \\
        20 M XOF & 1 993 & 1,29\,\% & 55,1\,\% & 2 063 XOF \\
        30 M XOF & 1 149 & 0,74\,\% & 48,0\,\% & 24 013 XOF \\
        40 M XOF & 764 & 0,49\,\% & 43,3\,\% & 41 276 XOF \\
        \textbf{50 M XOF} & \textbf{556} & \textbf{0,36\,\%} & \textbf{40,1\,\%} & \textbf{87 369 XOF} \\
        60 M XOF & 420 & 0,27\,\% & 37,5\,\% & 268 992 XOF \\
        75 M XOF & 304 & 0,20\,\% & 34,8\,\% & 493 012 XOF \\
        100 M XOF & 201 & 0,13\,\% & 31,7\,\% & 132 188 XOF \\
        \bottomrule
    \end{tabular}
\end{table}

À tous les seuils testés sur ces données, cet écart est positif, ce qui indique, sur les données étudiées, une séparation nette entre les deux populations plutôt qu'un chevauchement continu. Le seuil de 50 millions de XOF a été retenu comme compromis entre l'écart de séparation obtenu, la part d'encours qu'il permet d'isoler (40,1~\%) et le nombre de clients regroupés dans chaque segment.

Cette segmentation reste une règle de simulation, configurable, et non une classification juridique des clients : le fichier source ne contient aucune variable permettant d'identifier le statut juridique d'un client, et ce seuil sert de proxy opérationnel.

\subsection{Construction des variables explicatives et prévention des fuites de données}

Les variables explicatives sont construites séparément pour chaque segment (respectivement dans \texttt{retail\_features.py} et \texttt{corporate\_features.py}), selon trois catégories communes : des variables statiques, connues au moment de l'octroi ; des variables comportementales pour le segment Particuliers, ou de migration de rating pour le segment Entreprises, qui retracent l'historique du client jusqu'à la date d'observation ; et des variables macroéconomiques, propres à cette même date.

La construction de ces variables applique trois règles destinées à éviter toute fuite d'information entre les données disponibles au moment de la prédiction et celles utilisées pour l'évaluer. La variable de probabilité de défaut initiale est systématiquement exclue des variables explicatives. Le défaut à prédire, détaillé en section~3.3, est calculé sur une fenêtre strictement postérieure à la date d'observation : douze mois pour le segment Particuliers, quatre trimestres pour le segment Entreprises. Enfin, les comptes déjà en défaut à la date d'observation sont exclus, leur label n'étant pas défini.

\subsection{Organisation du Feature Store pour la modélisation}

Les variables ainsi construites sont organisées en deux pipelines distincts, l'un pour le segment Particuliers, l'autre pour le segment Entreprises. Cette séparation tient à la nature même des deux segments : leurs temporalités d'observation, mensuelle pour l'un, trimestrielle pour l'autre, ainsi que leurs variables explicatives, ne sont pas les mêmes, ce qui rend un traitement commun inapproprié.

Le choix de ces deux fréquences répond également à la nature des informations suivies. Pour les Particuliers, l'observation mensuelle permet de suivre les retards de paiement au rythme des échéances des crédits amortissables. Pour les Entreprises, le suivi repose davantage sur des ratios financiers et sur l'évolution de la notation, construits à partir d'informations comptables et financières dont l'actualisation n'est pas nécessairement mensuelle. Une granularité trimestrielle a donc été retenue dans le générateur afin de représenter cette évolution à un rythme adapté aux variables mobilisées. Ces fréquences constituent ainsi des choix de modélisation propres aux deux segments, et non des paramètres issus d'une comparaison expérimentale entre plusieurs cadences.

Les modèles construits à partir de ces deux pipelines sont présentés en section~3.4.

\section{Construction de la variable cible}

\subsection{Définition de l'événement de défaut}

Le défaut d'un client est défini différemment selon le segment, tout en reposant sur une logique commune : la confirmation de la dégradation sur plusieurs périodes consécutives, plutôt que sur un seul constat ponctuel. Pour le segment Particuliers, un client est considéré en défaut après trois mois consécutifs de retard de paiement sévère, de 90 jours ou plus. Pour le segment Entreprises, dont le suivi est trimestriel, un client est considéré en défaut après deux trimestres consécutifs dans la catégorie de notation CCC\_D, qui représente le niveau le plus dégradé de l'échelle retenue.

\subsection{Génération des historiques synthétiques}

En l'absence d'historique de crédit réel, rappelée en section~3.1, les trajectoires de remboursement et de défaut des clients ont été générées par simulation, à l'aide de chaînes de Markov : chaque client évolue, période après période, entre plusieurs états représentant son niveau de retard de paiement pour les Particuliers, ou sa notation pour les Entreprises, selon des probabilités de transition calibrées pour chaque segment. Ces transitions sont rendues sensibles à des variables macroéconomiques, le chômage et la croissance pour le segment Particuliers, la croissance et des taux propres à chaque secteur d'activité pour le segment Entreprises, de sorte que la probabilité de dégradation d'un client varie avec le contexte économique simulé plutôt que de rester constante dans le temps. Le détail des états et des matrices de transition retenues est présenté en annexe~C.

Les probabilités de transition sont calibrées autour de taux de défaut cibles de 7~\% pour le segment Particuliers et de 9~\% pour le segment Entreprises, à l'échelle des populations simulées. Ces taux cibles de calibration ne doivent pas être confondus avec le taux de défaut mesuré dans les échantillons d'évaluation des modèles, présenté en section~3.6, qui constitue une grandeur différente. Le taux cible du segment Entreprises a par ailleurs été révisé de 14~\% à 9~\%, après qu'une vérification du générateur a révélé une anomalie dans le mécanisme de simulation des notations. Après correction, le taux de 9~\% a été retenu comme valeur de calibration cohérente avec le comportement du générateur corrigé, et non comme un ajustement destiné à obtenir un résultat prédéfini.

\subsection{Construction des observations et horizons prédictifs}

Les trajectoires ainsi simulées sont ensuite observées à intervalles réguliers, un snapshot mensuel pour le segment Particuliers, trimestriel pour le segment Entreprises, afin de constituer les observations utilisées pour la modélisation. Chaque snapshot associe, à un client et à une date donnée, les variables explicatives construites en section~3.2 et l'indication de survenue ou non d'un défaut au cours de l'horizon prédictif suivant cette date : douze mois pour les Particuliers, quatre trimestres pour les Entreprises.

\section{Développement des modèles de PD}

\subsection{Démarche de modélisation}

Pour chaque segment, Particuliers et Entreprises, plusieurs algorithmes candidats sont entraînés et évalués selon un protocole commun, afin de retenir un modèle champion ; les autres algorithmes évalués sont conservés à titre de challengers. Cette architecture champion/challenger permet de documenter la comparaison entre plusieurs approches plutôt que de présenter un algorithme unique sans alternative examinée.

Le protocole appliqué à chaque algorithme candidat d'un même segment est identique : mêmes variables explicatives, mêmes jeux d'entraînement, de calibration et de test, et mêmes métriques d'évaluation, détaillées en section~3.6. Pour les modèles présentés dans cette section, qui répondent au besoin du suivi de portefeuille, le champion est désigné automatiquement parmi les candidats évalués, en retenant celui dont l'aire sous la courbe ROC (AUC), calculée sur le jeu de test, est la plus élevée. Le pipeline distinct propre à l'Origination, présenté en section~3.7, retient une démarche de sélection différente.

Cette méthode appelle une réserve méthodologique qu'il convient d'exposer clairement plutôt que de l'omettre. Le jeu de test sert à la fois à sélectionner le champion et à rapporter les performances finales de ce même champion, présentées en section~3.6. Cette double utilisation introduit un risque d'optimisme dans l'estimation de ces performances, puisque le modèle retenu l'est précisément parce qu'il obtient le meilleur résultat sur ce jeu de test. Ce risque n'est pas quantifié dans le cadre de ce projet. Il doit néanmoins être pris en compte dans l'interprétation des résultats de la section~3.6. La mise en place d'un jeu de validation distinct, réservé à la sélection du champion et laissant le jeu de test exclusivement consacré à l'évaluation finale, constitue à ce titre une amélioration méthodologique identifiée, mais non mise en \oe{}uvre dans le cadre de ce projet.

\subsection{Sélection des algorithmes et stratégie de validation par segment}

Les algorithmes candidats diffèrent entre les deux segments. Pour le segment Particuliers, une régression logistique et un modèle XGBoost sont comparés ; un modèle LightGBM a également été évalué, puis écarté avant la désignation du champion, pour un résultat insuffisant (AUC de 0,49, proche du hasard) et un surapprentissage constaté sur le découpage retenu pour ce segment. Pour le segment Entreprises, un modèle XGBoost, un modèle LightGBM et une régression logistique, retenue comme modèle de référence, sont comparés.

Les deux segments ne partagent ni la même temporalité d'observation ni la même structure d'identifiant de compte, rappelées en section~3.2 ; la stratégie de séparation entre entraînement et test diffère donc entre eux.

Pour le segment Particuliers, les comptes (et non directement les lignes) sont d'abord partitionnés en trois pools disjoints, entraînement, calibration et test, selon des proportions de 55~\%, 15~\% et 30~\%, stratifiées sur le statut \og jamais en défaut / au moins un défaut observé \fg{} afin de préserver le taux de défaut dans chaque pool ; l'absence de tout recouvrement de comptes entre ces pools est vérifiée de façon bloquante. Chaque pool ne fournit ensuite que les lignes correspondant à sa maturité : les pools d'entraînement et de calibration conservent les snapshots d'ancienneté 6 et 12 mois, le pool de test conserve le snapshot d'ancienneté 24 mois. Cette double séparation, par compte puis par ancienneté, corrige une fuite de données identifiée sur les données réelles du projet : un découpage antérieur, fondé sur l'ancienneté seule, laissait la totalité des comptes du jeu de test déjà présents dans le jeu d'entraînement à une ancienneté antérieure, ce qui permettait à un modèle capable de mémoriser de reconnaître un compte plutôt que d'être réellement évalué sur sa capacité de généralisation, les variables statiques d'un compte étant identiques entre ses différents snapshots. Une limite résiduelle, assumée et non corrigée dans ce lot, subsiste malgré cette correction : si les comptes des trois pools restent, eux, strictement disjoints, environ 60~\% des lignes du jeu d'entraînement portent sur un mois calendaire compris dans la plage de mois couverte par le jeu de test. La séparation par compte est ainsi assurée, mais la séparation calendaire reste partielle, l'horizon de simulation de 60 mois ne permettant pas un découpage calendaire complet sans perdre l'essentiel du volume disponible.

Pour le segment Entreprises, la séparation des données est construite à l'aide de \texttt{StratifiedGroupKFold}, avec cinq plis et un regroupement par compte. Cette méthode garantit qu'un même compte ne figure pas simultanément dans les ensembles d'entraînement et de validation d'un même pli. La stratification contribue par ailleurs à préserver une répartition comparable des défauts entre les plis, ceux-ci restant minoritaires. Parmi les cinq séparations ainsi générées, seule celle correspondant au dernier pli est retenue pour constituer les ensembles finaux d'entraînement et de test ; les quatre autres ne sont pas utilisées dans la sélection du champion ni dans le calcul des métriques finales présentées en section~3.6.

Ces deux choix sont d'abord justifiés par la structure propre des données de chaque segment, et par la nécessité d'éviter toute fuite d'information entre les observations utilisées pour l'entraînement et celles utilisées pour l'évaluation. Ils rejoignent par ailleurs, à titre de principes de validation plutôt que de prescription directe, certaines exigences formulées par le Comité de Bâle dans le cadre de l'approche notations internes (IRB), notamment les exigences relatives aux tests de performance hors période et hors échantillon (CRE36.36), ainsi qu'à la représentativité des données utilisées pour l'estimation (CRE36.67) \parencite{BCBS_CRE36}. Ces exigences réglementaires concernent les établissements qui utilisent l'approche IRB pour le calcul de leurs exigences de fonds propres ; elles ne s'appliquent pas en tant que telles à ce projet, qui ne constitue pas une validation réglementaire. Elles fournissent néanmoins un cadre de référence pour comprendre l'importance de la représentativité des données et de la distinction entre les observations utilisées pour construire un modèle et celles utilisées pour en évaluer les performances.

Le tableau~\ref{tab:strategies-validation} résume ces deux stratégies.

\begin{table}[htbp]
    \centering
    \footnotesize
    \caption{Synthèse des stratégies de validation par segment.}
    \label{tab:strategies-validation}
    \begin{tabular}{p{2.6cm}p{5.2cm}p{5.2cm}}
        \toprule
        \textbf{Élément} & \textbf{Particuliers} & \textbf{Entreprises} \\
        \midrule
        Stratégie & Partition des comptes en 3 pools disjoints, puis sélection par maturité & Validation croisée stratifiée et groupée par compte \\
        \addlinespace
        Méthode & Pools entraînement/calibration/test de comptes (55/15/30) ; entraînement et calibration : snapshots d'ancienneté 6 et 12 mois ; test : snapshot d'ancienneté 24 mois & \texttt{StratifiedGroupKFold}, 5 plis, groupe = compte ; seul le dernier pli retenu comme séparation finale \\
        \addlinespace
        Critère principal & Absence de recouvrement de comptes entre pools, puis ancienneté des observations & Regroupement par compte \\
        \addlinespace
        Objectif & Éviter qu'un compte déjà vu à l'entraînement ne soit reconnu plutôt qu'évalué sur sa généralisation, et observer le maintien des performances à une maturité plus élevée & Éviter la présence d'un même compte à la fois dans l'entraînement et dans la validation \\
        \addlinespace
        Précaution & Vérification bloquante de l'absence de recouvrement de comptes entre pools & Préservation de la répartition des défauts entre les plis \\
        \bottomrule
    \end{tabular}
\end{table}

\subsection{Entraînement et optimisation}

Les hyperparamètres des modèles XGBoost et LightGBM sont recherchés par optimisation bayésienne, à l'aide de la bibliothèque Optuna et de son échantillonneur TPE (\emph{Tree-structured Parzen Estimator}). Pour chaque segment, la fonction objectif maximisée est l'AUC moyenne obtenue par une validation croisée interne à 3 plis, réalisée uniquement sur le jeu d'entraînement : le jeu de test, décrit ci-dessus, n'intervient à aucun moment de cette recherche d'hyperparamètres.

Ce choix méthodologique est adapté à un espace de recherche combinant des hyperparamètres continus et discrets. Une recherche exhaustive par grille impose d'évaluer des combinaisons prédéfinies, tandis qu'une recherche aléatoire explore des configurations tirées indépendamment, sans exploiter directement les résultats des essais précédents. L'échantillonneur TPE adopte au contraire une démarche séquentielle : il modélise séparément les configurations associées aux meilleurs résultats et les autres, puis utilise ces estimations pour orienter les essais suivants vers des régions prometteuses de l'espace de recherche \parencite{BergstraEtAl2011}. Cette approche présente un intérêt lorsque chaque évaluation implique un nouvel entraînement du modèle et une validation croisée, dont le coût de calcul peut être important. Elle vise ainsi à exploiter plus efficacement le budget d'essais disponible, sans présumer d'une supériorité empiriquement démontrée dans le cadre de ce projet.

Cette validation croisée interne est une validation stratifiée simple à 3 plis, identique pour les deux segments ; elle n'est pas regroupée par compte pour le segment Entreprises, contrairement au découpage externe présenté précédemment. Un même compte peut donc apparaître à la fois dans l'ensemble d'apprentissage et dans celui de validation d'un pli au cours de cette recherche d'hyperparamètres. Cette situation introduit un risque de fuite d'information au sein de l'optimisation, susceptible de rendre les scores de validation interne plus favorables qu'ils ne le seraient avec une séparation stricte par compte. Cette limite concerne l'étape de recherche des hyperparamètres et doit être distinguée de la séparation externe, qui reste regroupée par compte.

\subsection{Calibration des probabilités}

Au-delà de leur capacité à ordonner correctement les emprunteurs selon leur risque, mesurée par l'AUC, les probabilités produites par les modèles doivent également être correctement calibrées, c'est-à-dire refléter un niveau de risque numériquement interprétable. Cette calibration, réalisée par une méthode de Platt ou une régression isotonique selon le cas, repose dans les deux segments sur des prédictions hors échantillon par rapport au modèle de PD qu'elle calibre, bien que ce caractère hors échantillon soit obtenu par deux mécanismes différents. Pour le segment Particuliers, la calibration utilise le pool de comptes dédié présenté en section~3.4.2, disjoint à la fois des comptes d'entraînement et des comptes de test. Pour le segment Entreprises, elle repose sur une procédure de validation croisée dédiée, décrite ci-après. Le jeu de test n'intervient pas dans l'ajustement des fonctions de calibration. Il reste toutefois utilisé pour la sélection du modèle champion et l'évaluation finale de ses performances, comme précisé en section~3.4.1.

Pour le segment Entreprises, la première stratégie de calibration reposait sur un échantillon qui ne contenait qu'une vingtaine de défauts, un nombre jugé insuffisant pour obtenir une calibration suffisamment fiable. Une calibration dite \emph{cross-fittée} a été mise en place pour remédier à cette limite : le modèle est entraîné sur plusieurs sous-ensembles du jeu d'entraînement, puis utilisé pour produire des prédictions hors échantillon sur la partie qu'il n'a pas vue à chacune de ces étapes ; l'ensemble de ces prédictions hors échantillon, une fois rassemblées, constitue l'échantillon de calibration effectivement utilisé. Cette méthode a permis de porter le nombre de défauts disponibles pour l'ajustement de la fonction de calibration d'une vingtaine à environ 135, avec des prédictions toutes produites hors échantillon : chaque observation y est évaluée par un modèle qui n'a pas été entraîné sur celle-ci.

\section{Développement du modèle de LGD}

\subsection{Construction de la variable cible et préparation des données}

Au-delà de la probabilité de défaut, la perte en cas de défaut (LGD) constitue la seconde composante de la perte attendue présentée au chapitre 1 : elle exprime, pour un compte en défaut, la part de l'exposition qui ne serait pas recouvrée. Comme pour la variable cible de la PD, cette grandeur n'est observable dans aucune donnée réelle du projet, rappelé en section~3.1 ; elle est construite par simulation, à partir d'un tirage aléatoire de type Beta du taux de recouvrement, conditionné par le type de garantie du compte. Ce taux brut est ensuite actualisé au taux directeur en vigueur au moment du défaut et net d'un coût de recouvrement (\emph{workout cost}), de sorte que la LGD obtenue, $LGD = 1 - \text{Recouvrement net}$, appartient à l'intervalle $[0,1]$ par construction, sans recours à un écrêtage artificiel.

Cette construction ne s'applique qu'aux comptes ayant effectivement connu un premier événement de défaut au cours de la trajectoire simulée, rappelée en section~3.3 : la LGD n'étant observable que pour ces comptes, la population d'entraînement est nécessairement restreinte à eux, soit 2~157 défauts pour le segment Particuliers et 135 défauts pour le segment Entreprises. Les variables explicatives retenues sont au nombre de deux : le ratio de couverture du crédit par sa garantie (\texttt{ratio\_couverture}) et le type de garantie (\texttt{type\_garantie}), cette dernière étant encodée par indicatrices (\emph{one-hot}).

\subsection{Développement et sélection du modèle}

Deux modèles candidats sont comparés pour prédire cette LGD. Une régression linéaire constitue le modèle de référence, interprétable mais susceptible de produire des valeurs hors de l'intervalle $[0,1]$ attendu pour une proportion. Une régression Beta, estimée par maximum de vraisemblance, lui est opposée comme modèle candidat principal : la LGD y est modélisée directement comme une variable suivant une loi Beta, $LGD \sim \text{Beta}(\mu\varphi, (1-\mu)\varphi)$, où $\mu$ est l'espérance prédite par le modèle (via une fonction sigmoïde des variables explicatives) et $\varphi$ un paramètre de précision estimé conjointement. Cette distribution est adaptée aux proportions bornées et permet de représenter une concentration aux extrémités de l'intervalle, par exemple une LGD proche de 0 pour un crédit bien garanti et proche de 1 pour un crédit sans garantie.

Le choix entre les deux modèles repose sur le critère d'information bayésien (BIC), qui tient compte de l'ajustement du modèle aux données tout en pénalisant sa complexité. Ce critère est utilisé ici pour comparer les deux approches sur les données d'entraînement ; il ne constitue toutefois pas une mesure de leur capacité de généralisation sur de nouvelles observations. Pour les deux segments, la régression Beta obtient un BIC inférieur à celui de la régression linéaire et est désignée championne : $-8\,175$ contre $-2\,534$ pour les Particuliers, $-242$ contre $-93$ pour les Entreprises.

\subsection{Prédiction et utilisation de la LGD}

Si l'entraînement du modèle de LGD est restreint, par nécessité, aux seuls comptes ayant connu un défaut, sa prédiction est en revanche appliquée à l'ensemble des comptes du portefeuille, qu'ils soient actuellement en défaut ou non : le calcul des pertes attendues présenté au chapitre 4 nécessite une LGD estimée pour chaque compte, à partir de ses seules caractéristiques de garantie, indépendamment de son état de défaut courant. La LGD ainsi prédite, combinée à la probabilité de défaut présentée en section~3.4 et à l'exposition au moment du défaut, alimente directement le calcul de la perte attendue développé au chapitre suivant.

\section{Évaluation, sélection et interprétabilité}

\subsection{Évaluation des performances et sélection des modèles}

La capacité des modèles à distinguer les emprunteurs selon leur niveau de risque est mesurée par trois indicateurs usuels en scoring de crédit. L'AUC (\emph{Area Under the Curve}) correspond à la probabilité qu'une paire d'emprunteurs, l'un en défaut et l'autre non, soit correctement ordonnée par le modèle ; elle varie entre 0,5 (pouvoir discriminant nul) et 1 (séparation parfaite). Le coefficient de Gini, défini par $2 \times \text{AUC} - 1$, en est une transformation linéaire d'usage courant dans le secteur bancaire. La statistique de Kolmogorov-Smirnov (KS) mesure l'écart maximal entre les distributions cumulées des scores des emprunteurs en défaut et de ceux qui ne le sont pas, et reflète la séparation la plus nette que le modèle parvient à établir entre les deux populations.

Le tableau~\ref{tab:performance-pd} présente ces trois indicateurs, calculés sur le jeu de test, pour l'ensemble des algorithmes candidats des deux segments.

\begin{table}[htbp]
    \centering
    \footnotesize
    \caption{Performances discriminantes des modèles de PD, par segment.}
    \label{tab:performance-pd}
    \begin{tabular}{p{2.4cm}p{4cm}p{1.8cm}p{1.8cm}p{1.8cm}}
        \toprule
        \textbf{Segment} & \textbf{Modèle} & \textbf{AUC} & \textbf{Gini} & \textbf{KS} \\
        \midrule
        Particuliers & Régression logistique (champion) & 0,7219 & 0,4438 & 0,3801 \\
        Particuliers & XGBoost & 0,6541 & 0,3082 & 0,2978 \\
        \addlinespace
        Entreprises & XGBoost (champion) & 0,8881 & 0,7761 & 0,6408 \\
        Entreprises & LightGBM & 0,8515 & 0,7031 & 0,5743 \\
        Entreprises & Régression logistique & 0,8122 & 0,6244 & 0,5227 \\
        \bottomrule
    \end{tabular}
\end{table}

Pour le segment Particuliers, la régression logistique domine nettement le modèle XGBoost sur les trois indicateurs, et est désignée championne conformément au critère décrit en section~3.4.1. Pour le segment Entreprises, XGBoost domine à la fois LightGBM et la régression logistique de référence, et est désigné champion selon le même critère ; l'écart entre les trois algorithmes y est toutefois plus resserré que pour les Particuliers. Ces désignations reposent, pour rappel, sur l'AUC du même jeu de test que celui utilisé pour rapporter les performances ci-dessus, ce qui introduit un risque d'optimisme déjà signalé en section~3.4.1 et qui doit être gardé à l'esprit dans la lecture de ces résultats, notamment pour le segment Entreprises où les écarts entre algorithmes sont les plus faibles.

\subsection{Évaluation de la calibration des probabilités}

La capacité d'un modèle à ordonner correctement les emprunteurs ne garantit pas que les probabilités qu'il produit soient numériquement fiables : un modèle peut bien classer les emprunteurs tout en leur attribuant des probabilités de défaut mal calibrées. Deux indicateurs complémentaires sont utilisés pour évaluer cette fiabilité. Le \emph{Brier Score} mesure l'erreur quadratique moyenne entre les probabilités prédites et les issues réellement observées (0 ou 1). Plus sa valeur est faible, meilleure est la qualité globale des probabilités produites. Cet indicateur ne mesure toutefois pas exclusivement la calibration, puisqu'il dépend également de la capacité discriminante du modèle. L'\emph{Expected Calibration Error} (ECE) regroupe les prédictions en dix classes de probabilité croissante et mesure l'écart moyen, pondéré par la taille de chaque classe, entre la probabilité moyenne prédite et la fréquence de défaut réellement observée dans cette classe.

Le tableau~\ref{tab:calibration-pd} présente ces deux indicateurs, ainsi que la probabilité moyenne calibrée (\texttt{pd\_moy\_cal}) et le taux de défaut mesuré dans l'échantillon de test (\texttt{taux\_reel}), pour les mêmes modèles.

\begin{table}[htbp]
    \centering
    \footnotesize
    \caption{Calibration des probabilités de PD, par segment.}
    \label{tab:calibration-pd}
    \begin{tabular}{p{2.4cm}p{3.6cm}p{1.6cm}p{1.6cm}p{2cm}p{1.8cm}}
        \toprule
        \textbf{Segment} & \textbf{Modèle} & \textbf{Brier} & \textbf{ECE} & \textbf{PD moyenne calibrée} & \textbf{Taux mesuré} \\
        \midrule
        Particuliers & Régression logistique (champion) & 0,0218 & 0,0093 & 0,0325 & 0,0232 \\
        Particuliers & XGBoost & 0,0218 & 0,0067 & 0,0273 & 0,0232 \\
        \addlinespace
        Entreprises & XGBoost (champion) & 0,0212 & 0,0044 & 0,0249 & 0,0256 \\
        Entreprises & LightGBM & 0,0222 & 0,0111 & 0,0228 & 0,0256 \\
        Entreprises & Régression logistique & 0,0226 & 0,0074 & 0,0233 & 0,0256 \\
        \bottomrule
    \end{tabular}
\end{table}

Ce taux mesuré dans l'échantillon de test (2,32~\% pour les Particuliers, 2,56~\% pour les Entreprises) ne doit pas être confondu avec les taux cibles de calibration du générateur synthétique, 7~\% et 9~\% respectivement, présentés en section~3.3 : ces deux grandeurs s'observent à des échelles différentes, celle de la population simulée dans son ensemble pour les taux cibles, celle du seul échantillon de test pour le taux mesuré ici. Pour les deux segments, la PD moyenne calibrée des champions reste proche du taux mesuré, sans s'y confondre exactement, ce qui est attendu d'une calibration opérant sur l'ensemble de la distribution des scores plutôt que sur sa seule moyenne. Il convient enfin de noter que les performances discriminantes et la qualité de la calibration ne se recoupent pas nécessairement : le champion Entreprises obtient à la fois la meilleure discrimination et la meilleure calibration parmi les candidats de son segment, mais ce n'est pas une propriété garantie par construction, un modèle pouvant en principe bien discriminer sans être bien calibré, ou l'inverse.

\subsection{Interprétabilité des modèles par SHAP}

Au-delà de leurs performances globales, les modèles retenus doivent fournir des éléments permettant de comprendre leurs résultats individuels, notamment dans le cadre de l'instruction d'une demande de crédit. Cette explicabilité repose sur les valeurs de Shapley (SHAP, \emph{SHapley Additive exPlanations}), issues de la théorie des jeux coopératifs \parencite{LundbergLee2017} : elles décomposent le score brut produit par un modèle pour une observation donnée en une contribution additive de chacune de ses variables explicatives, de sorte que la somme de ces contributions et d'une valeur de référence reconstitue exactement ce score. Le calcul de ces contributions utilise un explainer exact adapté à chaque algorithme : un \emph{TreeExplainer} pour les modèles XGBoost et LightGBM, un \emph{LinearExplainer} pour la régression logistique.

Deux niveaux d'explication sont effectivement produits. Au niveau global, l'importance moyenne absolue de chaque variable sur l'ensemble du portefeuille fournit une vue d'ensemble des facteurs les plus influents pour un modèle donné, utile à l'audit du modèle et à sa documentation. Au niveau individuel, les principaux facteurs de risque et facteurs protecteurs d'un compte particulier sont identifiés afin de fournir des éléments d'explication sur le score attribué par le modèle. Dans le parcours Origination, ces explications sont également intégrées à la trace d'audit associée à la décision, comme détaillé en section~3.7.5.

Ces explications portent sur les scores produits par le modèle avant calibration, la calibration étant appliquée séparément, comme décrit en section~3.4.4 ; elles décrivent donc la contribution de chaque variable au score du modèle, non aux probabilités calibrées finalement restituées à l'utilisateur. Deux limites doivent par ailleurs être gardées à l'esprit dans leur interprétation. D'une part, une contribution SHAP élevée indique l'importance d'une variable pour la prédiction d'\emph{un modèle donné} ; elle ne constitue pas une preuve que cette variable est, dans la réalité, une cause du défaut. D'autre part, puisque les modèles présentés dans ce chapitre ont été entraînés sur des historiques de défaut synthétiques, rappelé en section~3.1, les facteurs mis en évidence par SHAP reflètent la façon dont le générateur de données a construit la relation entre les variables et le défaut, et non nécessairement la dynamique réelle du risque de crédit dans un portefeuille bancaire.

\subsection{Évaluation des modèles de LGD}

L'évaluation des modèles de LGD, présentés en section~3.5, appelle une précaution méthodologique distincte de celle des modèles de PD. Le tableau~\ref{tab:performance-lgd} présente, pour les deux segments, l'erreur quadratique moyenne (RMSE), le coefficient de détermination ($R^2$) et le BIC des deux modèles candidats.

\begin{table}[htbp]
    \centering
    \footnotesize
    \caption{Performances des modèles de LGD, par segment.}
    \label{tab:performance-lgd}
    \begin{tabular}{p{2.4cm}p{3.2cm}p{1.8cm}p{1.8cm}p{2cm}}
        \toprule
        \textbf{Segment} & \textbf{Modèle} & \textbf{RMSE} & \textbf{$R^2$} & \textbf{BIC} \\
        \midrule
        Particuliers & Régression linéaire & 0,1328 & 0,8000 & $-2\,534$ \\
        Particuliers & Beta Regression (champion) & 0,1329 & 0,7998 & $-8\,175$ \\
        \addlinespace
        Entreprises & Régression linéaire & 0,1512 & 0,6426 & $-93$ \\
        Entreprises & Beta Regression (champion) & 0,1512 & 0,6424 & $-242$ \\
        \bottomrule
    \end{tabular}
\end{table}

Ces valeurs appellent une lecture différente de celles présentées pour la PD. D'une part, à la différence des modèles de PD, dont l'AUC/Gini/KS et les indicateurs de calibration sont calculés sur un jeu de test distinct de l'entraînement, les métriques de LGD ci-dessus sont calculées sur les mêmes données que celles utilisées pour entraîner les modèles : elles ne constituent pas une évaluation hors échantillon et ne doivent pas être lues comme telle. Le projet documente d'ailleurs lui-même cette distinction, en qualifiant explicitement cette base d'évaluation d'\emph{in-sample} pour la LGD, par opposition à la validation hors échantillon réelle dont bénéficient les modèles de PD. D'autre part, l'échantillon d'entraînement du segment Entreprises ne compte que 135 défauts, un volume restreint qui limite la robustesse des conclusions qui peuvent en être tirées, au-delà même de la question de l'échantillon d'évaluation. Le RMSE et le $R^2$ très proches entre les deux modèles candidats, pour les deux segments, confirment que l'écart entre eux tient avant tout à la pénalisation de complexité du BIC plutôt qu'à une différence marquée de leur ajustement aux données.

\section{Modélisation prédictive spécifique au parcours Origination}

\subsection{Spécificités et objectifs du modèle PD à l'Origination}

Les modèles de probabilité de défaut présentés en section~3.4 répondent au besoin du suivi de portefeuille : réévaluer, à intervalles réguliers, le risque de comptes déjà en portefeuille, pour lesquels un historique de remboursement s'est progressivement constitué. L'évaluation d'une nouvelle demande de crédit à l'Origination pose un problème différent : le crédit demandé n'a pas encore été octroyé, et les informations relatives à son propre comportement futur, retard de paiement, dégradation de notation, utilisation d'une ligne de trésorerie, n'existent donc pas au moment de la décision. Certaines informations comportementales antérieures concernant le client, par exemple issues de ses autres comptes, pourraient en revanche théoriquement être exploitées ; leur disponibilité effective dépend alors de ce que la plateforme collecte réellement, une distinction précisée plus loin dans cette sous-section. C'est cette contrainte qui explique la distinction, déjà présentée au chapitre 5, entre les modèles de probabilité de défaut utilisés par les deux parcours : la plateforme entraîne, par un processus distinct, deux modèles de PD propres à l'Origination, un pour chaque segment, tout en réutilisant les composantes de LGD, d'EAD et de rentabilité déjà mobilisées dans les deux parcours, dont les mécanismes sont détaillés aux chapitres 3 et 4.

Cette contrainte de disponibilité se traduit concrètement dans les variables explicatives retenues. Les deux ensembles de variables propres à l'Origination sont construits, dans le code du projet, comme des sous-ensembles stricts des variables utilisées pour le suivi de portefeuille, une relation vérifiée automatiquement plutôt que simplement documentée : toute variable comportementale, construite à partir d'un historique postérieur à l'octroi, en est exclue. Un cas précis illustre le soin apporté à cette contrainte : deux variables, le nombre de restructurations antérieures et le pire retard de paiement déjà observé, constituent bien, en toute rigueur, des informations potentiellement disponibles avant l'octroi d'un nouveau crédit ; elles ont néanmoins été exclues du modèle d'Origination, comme le précise le code, au motif qu'aucun canal de collecte réel, formulaire ou interface de saisie, ne permet aujourd'hui de les renseigner effectivement au moment d'une demande. Cette exclusion ne répond donc pas à une indisponibilité théorique de l'information, mais à l'absence de collecte effective par la plateforme telle qu'elle est réalisée aujourd'hui, un choix propre à l'état actuel de la plateforme plutôt qu'une contrainte méthodologique générale.

\subsection{Constitution de la cible et découpage temporel}

La variable cible du modèle d'Origination n'introduit pas une nouvelle définition du défaut : elle reprend exactement la même définition et le même horizon de douze mois que la variable cible du suivi de portefeuille, présentée en section~3.3, en modifiant uniquement le point d'ancrage de cet horizon. Pour le suivi de portefeuille, l'horizon de prédiction est calculé à partir de chaque date de snapshot ; pour l'Origination, il est ancré sur la période d'octroi du crédit, le mois pour les Particuliers et le trimestre pour les Entreprises. Il ne s'agit donc pas d'une règle métier nouvelle, mais d'un changement mécanique de ce point de départ. Un audit de faisabilité dédié a notamment vérifié le mécanisme de sauvegarde contre le biais de survie : pour le segment Particuliers, aucun compte éligible n'a été exclu par ce mécanisme en raison d'un historique insuffisant pour couvrir l'horizon de prédiction.

La séparation entre les données d'entraînement et de test ne reprend pas davantage le découpage par pools de comptes présenté en section~3.4.2. Elle repose sur un découpage chronologique selon la date d'Origination : pour les Particuliers, les comptes sont ordonnés par mois d'octroi ; pour les Entreprises, ils le sont par trimestre. Dans chaque segment, les périodes les plus anciennes sont affectées à l'entraînement et les plus récentes au test, selon une proportion de 80~\% contre 20~\%. Cette stratégie vise à évaluer le modèle sur des octrois plus récents que ceux utilisés pour son entraînement. Cette proportion constitue un paramètre retenu spécifiquement pour le pipeline Origination : elle ne reprend pas une valeur préexistante déjà formalisée pour le suivi de portefeuille.

\subsection{Entraînement, optimisation et sélection du modèle}

Les candidats XGBoost et LightGBM du pipeline Origination sont entraînés selon le même mécanisme que celui présenté en section~3.4.3 pour le suivi de portefeuille : la même classe du projet est réutilisée, et leurs hyperparamètres sont recherchés par la même optimisation bayésienne Optuna/TPE, sur une validation croisée interne limitée au jeu d'entraînement. Cette classe désigne également, par défaut et automatiquement, un champion parmi ses candidats en retenant celui dont l'AUC sur le jeu de test est la plus élevée, exactement comme pour le suivi de portefeuille. Ce comportement par défaut n'a pas été modifié : il reste identique pour l'ensemble des usages de la classe, y compris les modèles de suivi existants. Pour l'Origination cependant, cette désignation automatique n'est pas celle qui déterminera le champion final : elle est explicitement écartée au profit d'une décision prise en dehors de la classe, à l'issue d'un protocole de validation dédié.

Ce protocole, exécuté et validé par l'équipe projet, répond à une règle énoncée explicitement dans sa documentation : aucun modèle ne doit entrer dans le moteur de décision sur la seule base de son AUC, avant d'avoir satisfait un protocole de validation du risque complet, examinant successivement l'absence de fuite de données, la stabilité temporelle, la calibration, la cohérence des facteurs explicatifs SHAP avec la logique métier attendue, puis une comparaison multicritère conduisant à la décision de champion. Le tableau~\ref{tab:comparaison-origination} résume les quatre critères ainsi comparés, discrimination, calibration, stabilité temporelle et interprétabilité, pour les candidats des deux segments.

\begin{table}[htbp]
    \centering
    \footnotesize
    \caption{Comparaison multicritère des candidats PD Origination, par segment.}
    \label{tab:comparaison-origination}
    \begin{tabular}{p{1.8cm}p{2.2cm}p{1.6cm}p{2.6cm}p{2.8cm}p{2.6cm}}
        \toprule
        \textbf{Segment} & \textbf{Modèle} & \textbf{AUC} & \textbf{Pente de calibration (cible = 1)} & \textbf{Écart de stabilité temporelle} & \textbf{Interprétabilité} \\
        \midrule
        Particuliers & Régression logistique (champion) & 0,6546 & 1,012 & 9,5 points & Interprétation directe des coefficients \\
        Particuliers & XGBoost & 0,6544 & 1,018 & 13,5 points & Explication complémentaire par SHAP \\
        \addlinespace
        Entreprises & Régression logistique (champion) & 0,772 & 0,715 & 11,7 points & Interprétation directe des coefficients \\
        Entreprises & XGBoost & 0,762 & 0,516 & 16,1 points & Explication complémentaire par SHAP \\
        Entreprises & LightGBM & 0,778 & 0,420 & 16,6 points & Explication complémentaire par SHAP \\
        \bottomrule
    \end{tabular}
\end{table}

L'écart de stabilité temporelle correspond à la différence, exprimée en points d'AUC, entre les performances mesurées sur les deux sous-périodes comparées du jeu de test ; une valeur plus faible traduit une variation moins importante des performances entre ces deux sous-périodes.

Ces résultats ne conduisent pas à la même lecture pour les deux segments, et il serait inexact de présenter la régression logistique comme supérieure aux autres candidats sur l'ensemble des critères. Pour le segment Particuliers, l'AUC de la régression logistique et celle du modèle XGBoost sont quasiment identiques, et leur calibration est elle aussi très proche, les deux modèles restant proches de la pente idéale ; la calibration ne départage donc pas les deux candidats. Ce sont la stabilité temporelle, meilleure pour la régression logistique, et son interprétabilité directe, qui motivent le choix du champion sur ce segment. Pour le segment Entreprises, c'est LightGBM qui obtient l'AUC la plus élevée des trois candidats, ce qui favoriserait, sur ce seul critère, sa désignation comme champion. Cet avantage n'est toutefois pas jugé suffisamment établi : le jeu de test Entreprises ne comporte que 227 lignes pour 13 défauts, et une estimation de la variabilité de l'AUC à cette taille d'échantillon, par une procédure de rééchantillonnage, situe l'écart observé entre LightGBM et la régression logistique dans la marge d'incertitude attendue à cette échelle, plutôt que comme un écart de discrimination réellement démontré. C'est alors la calibration, nettement plus favorable à la régression logistique, qui devient le critère déterminant pour ce segment.

La décision retenue à l'issue de ce protocole est ainsi la régression logistique comme champion pour les deux segments du pipeline Origination, une décision méthodologique distincte de celle retenue pour le suivi de portefeuille, où la régression logistique n'est championne que pour le segment Particuliers, XGBoost l'étant pour le segment Entreprises, comme présenté en section~3.6.1. Ces deux décisions, prises sur deux jeux de données et selon deux protocoles différents, ne doivent pas être confondues.

\subsection{Calibration et seuils opérationnels}

La calibration des probabilités d'Origination repose sur le même principe que celui présenté en section~3.4.4 pour le suivi de portefeuille : les fonctions de calibration sont ajustées sur des prédictions hors échantillon par rapport au modèle de PD qu'elles calibrent. Le mécanisme diffère néanmoins par segment. Pour les Particuliers, aucun pool de comptes dédié à la calibration n'est construit pour l'Origination, à la différence du suivi de portefeuille ; c'est alors le comportement par défaut de la classe qui s'applique, les 20~\% les plus récents du jeu d'entraînement, au sein de la période d'entraînement elle-même et toujours antérieurs à la période de test, servant d'échantillon de calibration. Pour les Entreprises, la même calibration cross-fittée que celle présentée pour le suivi de portefeuille est réutilisée, pour le même motif, un nombre de défauts trop restreint pour ajuster un calibrateur fiable à partir d'un seul jeu de calibration ; la procédure est toutefois plus simple ici, une validation croisée stratifiée ordinaire suffisant, puisque le jeu d'Origination compte une ligne par compte et non plusieurs lignes de panel par compte comme au suivi de portefeuille.

Indépendamment de cette calibration, un second mécanisme, distinct de la politique de décision présentée en section~4.5.2, détermine un seuil de probabilité de défaut par segment à partir d'un compromis entre les coûts estimés des erreurs de classification. Ce compromis est traduit par un score F-bêta, dont le paramètre $\beta$ pondère le coût d'un faux négatif, approximé par le produit de la LGD moyenne et de l'exposition moyenne au moment du défaut, et celui d'un faux positif, approximé par un taux de marge appliqué à cette même exposition moyenne. Ce seuil est recalculé spécifiquement sur les données de test du pipeline Origination, et non recopié depuis le modèle de suivi, ce qui conduit à des valeurs numériques distinctes, de l'ordre de 3,3~\% pour les Particuliers et de 2,8~\% pour les Entreprises. Le jeu de test intervient donc ici dans l'ajustement du seuil d'alerte, contrairement à son rôle dans l'évaluation des performances discriminantes du modèle présentée en section~3.7.3 ; les valeurs obtenues doivent être considérées comme des paramètres empiriques du dispositif d'alerte sur l'échantillon étudié, et non comme des seuils dont la performance aurait été confirmée sur un échantillon indépendant.

Ce seuil n'intervient à aucun moment dans la décision d'octroi elle-même, qui reste entièrement gouvernée par la politique de décision présentée en section~4.5.2. Il alimente uniquement un signal distinct, calculé au moment même de la décision pour les seuls comptes approuvés : ceux dont la probabilité de défaut dépasse ce seuil sont signalés, dans la trace de la décision elle-même, comme justifiant une vigilance renforcée. Ce signal reste descriptif : il n'est pas conservé comme un état de suivi autonome et est réévalué à chaque décision plutôt que déclenché une seule fois puis mémorisé. Le seuil F-bêta constitue ainsi un paramètre de ce signal, et non un critère supplémentaire d'acceptation ou de refus du crédit, ni le déclencheur d'un processus de surveillance distinct et persistant.

\subsection{Explicabilité SHAP du modèle d'Origination}

La cohérence des facteurs explicatifs SHAP avec la logique métier attendue a déjà été examinée en section~3.7.3, dans le cadre du protocole de validation ayant conduit à la désignation de la régression logistique comme champion pour les deux segments. Cet examen portait sur les candidats eux-mêmes, à titre diagnostic, et ne suppose pas qu'un mécanisme d'explicabilité soit effectivement disponible lors de l'évaluation d'une nouvelle demande. Un second chantier, distinct et postérieur, a porté sur la construction d'un tel mécanisme et son intégration effective au parcours Origination.

Puisque la régression logistique est championne pour les deux segments du pipeline Origination, l'explicabilité repose sur un \emph{LinearExplainer}, celui des deux explainers présentés en section~3.6.3 qui correspond à cet algorithme. Un \emph{SHAPEngine} dédié à l'Origination, distinct de celui du suivi de portefeuille, est ajusté séparément pour chaque segment, sur un échantillon de comptes construit à l'aide des mêmes fonctions de préparation des données et de construction des variables que celles utilisées pour l'entraînement des modèles d'Origination, et non à partir d'une approximation distincte. Les modèles de PD eux-mêmes ne sont pas réentraînés à cette occasion.

Ce mécanisme a été intégré à l'artefact de décision Origination par une procédure prudente : une copie de sauvegarde de cet artefact est conservée avant toute modification, l'écriture s'effectue d'abord sur un fichier temporaire, et ce fichier temporaire est rechargé et vérifié de façon indépendante, notamment en exécutant effectivement une décision complète sur des comptes réels des deux segments et en confirmant qu'une explication est bien produite, avant que l'artefact initial ne soit remplacé. Cette précaution tient au fait que la procédure modifie un artefact déjà en usage plutôt que d'en produire un nouveau : un échec à n'importe quelle étape laisse l'artefact initial inchangé.

\section*{Conclusion du chapitre}

Ce chapitre a retracé l'ensemble de la démarche qui mène d'une source de données réelle, mais limitée, à des modèles de risque de crédit évalués et, pour la PD, interprétables. La source réelle, un fichier de comptes de dépôt, a fourni une structure de clientèle authentique, mais ne contenait ni historique de crédit ni défaut observé ; cette limite a rendu nécessaire la génération d'historiques de défaut et de recouvrement synthétiques, construits par simulation et calibrés par segment, dont la construction a été détaillée et dont les limites ont été assumées plutôt que dissimulées. Sur ces données, deux composantes de la perte attendue ont été modélisées. Pour la probabilité de défaut, deux pipelines distincts ont été développés. Dans le parcours Portfolio Analytics, adapté aux caractéristiques propres des segments Particuliers et Entreprises, et avec une attention portée à la prévention des fuites de données, y compris lorsqu'une telle fuite a été identifiée et corrigée en cours de projet, la régression logistique a été retenue pour les Particuliers et XGBoost pour les Entreprises, selon une sélection fondée sur l'AUC du jeu de test. Dans le parcours Origination, une validation multicritère intégrant discrimination, calibration, stabilité temporelle et interprétabilité a conduit à retenir la régression logistique pour les deux segments. Ces modèles ont ensuite été évalués selon leur capacité de discrimination et la fiabilité de leurs probabilités, sur un jeu de test distinct de l'entraînement, puis rendus interprétables par les valeurs de Shapley. Pour la perte en cas de défaut, une régression Beta a été retenue pour les deux segments, entraînée sur la seule population des comptes en défaut et évaluée en conditions in-sample, une limite assumée plutôt que dissimulée. Les réserves méthodologiques rencontrées au long de ce cheminement, qu'elles concernent la sélection du champion de PD, l'évaluation in-sample de la LGD ou la nature synthétique des données d'entraînement, ont été exposées explicitement. Les performances présentées attestent ainsi du comportement des modèles sur les données construites pour le projet, sans démontrer leur capacité prédictive sur des défauts bancaires réels. Le chapitre suivant mobilise ces modèles de probabilité de défaut et de perte en cas de défaut, aux côtés d'une exposition au défaut calculée mécaniquement, pour aborder la quantification du risque, du capital et de la rentabilité, jusqu'à l'aide à la décision d'octroi.
